\chapter{ĐỒNG BỘ CẤU HÌNH}
\label{chapter:chapter-4}
% Yêu cầu package: \usepackage{tikz} và \usetikzlibrary{arrows.meta,shapes,fit}

% ============================================================
\section{Tổng Quan}
% ============================================================

Đồng bộ cấu hình là quá trình \textbf{Bộ điều khiển (Controller)} thiết lập và duy trì
các tham số vận hành thống nhất cho tất cả \textbf{Tác nhân (Agent)} trong mạng
Multi-AP~[1]. Thay vì mỗi Agent tự cấu hình độc lập, mô hình tập trung này đảm bảo mọi
điểm truy cập phát sóng cùng SSID, áp dụng cùng chính sách chuyển vùng và tuân thủ
cùng quy tắc phân tách lưu lượng.

Quá trình đồng bộ bao gồm ba lớp cấu hình chính:
\begin{itemize}
	\item \textbf{Cấu hình AP (AP Configuration):} Thiết lập các tham số radio và BSS
	(SSID, xác thực, vai trò Fronthaul/Backhaul) thông qua WSC hoặc DPP~[1, 2].
	\item \textbf{Đồng bộ Chính sách (Policy Sync):} Phân phối các quy tắc điều phối,
	ngưỡng chuyển vùng và ánh xạ VLAN tới từng Agent~[15--20].
	\item \textbf{Cấu hình MLO (Wi-Fi 7):} Quản lý các thực thể đa liên kết (MLD) và
	các Affiliated AP trong môi trường Wi-Fi 7~[21, 22].
\end{itemize}

Định danh radio trong tất cả các bản tin đồng bộ được thực hiện qua \textbf{RUID}
(\textit{Radio Unique Identifier}) --- một giá trị MAC 6 byte gán duy nhất cho mỗi radio
trên Agent, đảm bảo Controller luôn biết cấu hình đang áp dụng cho radio nào.


% ============================================================
\section{Cơ Chế Cấu Hình AP (AP Configuration)}
% ============================================================

Việc cấu hình các radio và BSS trên Agent được thực hiện qua hai giao thức tùy theo
thế hệ thiết bị: \textbf{WSC} (Wi-Fi Simple Configuration) dành cho Profile-1/2 và
\textbf{DPP} (Device Provisioning Protocol) dành cho Profile-3 trở lên.


% ------------------------------------------------------------
\subsection{Cấu Hình Qua WSC (Wi-Fi Simple Configuration)}
% ------------------------------------------------------------

WSC là phương thức cấu hình cơ bản, được mở rộng từ IEEE 1905.1 để hỗ trợ Multi-AP~[3, 4].
Quá trình diễn ra theo mô hình \textit{Registrar} (Controller) -- \textit{Enrollee} (Agent):
Agent chủ động tìm kiếm Controller, sau đó hai bên trao đổi khả năng và nhận cấu hình
theo từng radio.

\begin{itemize}
	\item \textbf{M1 Message:} Agent gửi thông tin năng lực của radio thông qua
	\textit{AP Radio Basic Capabilities TLV}, bao gồm các băng tần hỗ trợ, số
	chuỗi không gian (spatial streams) và các chuẩn Wi-Fi tương thích~[5, 6].
	\item \textbf{M2 Message:} Controller phản hồi với cấu hình BSS cụ thể --- SSID,
	phương thức xác thực (WPA2/WPA3) --- và chỉ định vai trò của BSS thông qua
	\textit{Multi-AP Extension subelement}: Fronthaul (kết nối thiết bị người dùng),
	Backhaul (kết nối giữa các Agent), hoặc cả hai~[7, 8].
\end{itemize}

\textbf{Lưu ý:} Quy trình M1/M2 được lặp lại độc lập cho \textit{từng radio} trên Agent.
Agent có bao nhiêu radio thì sẽ có bấy nhiêu cặp trao đổi M1/M2 riêng biệt.


% ------------------------------------------------------------
\subsection{Sơ Đồ Luồng Bản Tin WSC Configuration}
% ------------------------------------------------------------

\begin{center}
	\begin{tikzpicture}[>=stealth, font=\small]
		% --- Tiêu đề ---
		\node[draw, fill=gray!15, minimum width=2.8cm, minimum height=0.6cm]
		(C) at (0, 0) {Controller};
		\node[draw, fill=gray!15, minimum width=2.8cm, minimum height=0.6cm]
		(A) at (8, 0) {Multi-AP Agent};
		
		% --- Lifelines ---
		\draw[dashed, gray!50] (0, -0.4) -- (0, -10.5);
		\draw[dashed, gray!50] (8, -0.4) -- (8, -10.5);
		
		% --- Bước 1 & 2: Discovery ---
		\draw[->] (8,-1.2) -- (0,-1.2)
		node[midway, above, scale=0.7, fill=white, inner sep=1pt]
		{(1) 1905 AP-Autoconfig Search};
		\draw[->] (0,-2.2) -- (8,-2.2)
		node[midway, above, scale=0.7, fill=white, inner sep=1pt]
		{(2) 1905 AP-Autoconfig Response};
		
		% --- Nhãn giai đoạn ---
		\draw[dashed, gray!70] (-1.2,-3.0) rectangle (9.2,-7.8);
		\node[scale=0.65, gray!80, anchor=west] at (-1.15,-3.25)
		{Lặp lại cho từng Radio (xác định qua RUID)};
		
		% --- Bước 3 & 4: WSC M1/M2 ---
		\draw[->, blue] (8,-4.2) -- (0,-4.2)
		node[midway, above, scale=0.7, fill=white, inner sep=1pt]
		{(3) WSC M1 [AP Radio Basic Capabilities TLV]};
		\draw[->, blue] (0,-5.6) -- (8,-5.6)
		node[midway, above, scale=0.7, fill=white, inner sep=1pt]
		{(4) WSC M2 [BSS Config: SSID, AKM, Roles]};
		
		% --- Nhãn giải thích M2 ---
		\node[scale=0.6, align=left, anchor=west] at (0.3,-6.6)
		{\textit{Multi-AP Extension subelement:}
			Fronthaul / Backhaul / cả hai};
		
		% --- Bước 5: Topology Report ---
		\draw[->, brown] (8,-8.8) -- (0,-8.8)
		node[midway, above, scale=0.7, fill=white, inner sep=1pt]
		{(5) 1905 Topology Response [AP Operational BSS TLV]};
		\node[scale=0.6, gray, align=center] at (4,-9.5)
		{\textit{Agent báo cáo BSSID thực tế đang phát sóng để Controller cập nhật bản đồ mạng}};
	\end{tikzpicture}
\end{center}

\textbf{Giải thích luồng bản tin:}
\begin{enumerate}
	\item Agent phát \textit{AP-Autoconfig Search} trên mạng 1905.1 để tìm kiếm
	Controller đóng vai trò Registrar.
	\item Controller phản hồi, xác nhận vai trò và thiết lập phiên cấu hình.
	\item Với mỗi radio, Agent gửi \textit{WSC M1} kèm \textit{AP Radio Basic
		Capabilities TLV} mô tả chi tiết năng lực phần cứng (RUID, chuẩn Wi-Fi,
	số chuỗi không gian, băng tần).
	\item Controller phân tích khả năng radio và phản hồi \textit{WSC M2} với cấu hình
	BSS phù hợp, kèm \textit{Multi-AP Extension subelement} chỉ định rõ vai trò
	từng BSS.
	\item Sau khi áp dụng cấu hình, Agent gửi \textit{Topology Response} với
	\textit{AP Operational BSS TLV} liệt kê các BSSID thực tế đang hoạt động,
	giúp Controller cập nhật bản đồ mạng.
\end{enumerate}


% ============================================================
\section{Cấu Hình Qua DPP (Device Provisioning Protocol)}
% ============================================================

Đối với các thiết bị hỗ trợ \textbf{Profile-3}, cấu hình BSS được thực hiện sau khi
đã thiết lập kênh bảo mật 1905 (PMK/PTK) thông qua quy trình DPP Onboarding~[9, 10].
Khác với WSC dùng định dạng nhị phân, DPP sử dụng \textbf{đối tượng JSON} để mã hóa
cấu hình, cho phép biểu diễn linh hoạt hơn các tham số phức tạp của Wi-Fi 6/7.

\begin{itemize}
	\item \textbf{BSS Configuration Request:} Agent báo cáo chi tiết năng lực của từng
	radio theo các chuẩn HT (802.11n), VHT (802.11ac), HE (802.11ax) và EHT
	(802.11be/Wi-Fi 7), giúp Controller lựa chọn cấu hình tối ưu cho từng
	phần cứng~[10, 11].
	\item \textbf{BSS Configuration Response:} Controller cung cấp \textit{DPP Config
		Object} ở định dạng JSON cho từng BSS, bao gồm thông tin SSID, phương thức
	xác thực, thông số MLD (nếu có) và chứng chỉ mạng~[12--14].
\end{itemize}


% ------------------------------------------------------------
\subsection{Sơ Đồ Luồng Bản Tin DPP Configuration}
% ------------------------------------------------------------

\begin{center}
	\begin{tikzpicture}[>=stealth, font=\small]
		% --- Tiêu đề ---
		\node[draw, fill=gray!15, minimum width=2.8cm, minimum height=0.6cm]
		(C) at (0, 0) {Controller};
		\node[draw, fill=gray!15, minimum width=2.8cm, minimum height=0.6cm]
		(A) at (8, 0) {Multi-AP Agent};
		
		% --- Lifelines ---
		\draw[dashed, gray!50] (0,-0.4) -- (0,-10);
		\draw[dashed, gray!50] (8,-0.4) -- (8,-10);
		
		% --- Điều kiện tiên quyết ---
		\node[draw, dashed, fill=green!10, scale=0.72, align=center] at (4,-1.1)
		{Điều kiện: Đã thiết lập 1905 Secure Link (PMK/PTK)};
		
		% --- Bước 1: BSS Config Request ---
		\draw[->, orange] (8,-2.4) -- (0,-2.4)
		node[midway, above, scale=0.7, fill=white, inner sep=1pt]
		{(1) BSS Configuration Request [HT/VHT/HE/EHT Caps]};
		
		% --- Bước 2: BSS Config Response ---
		\draw[->, orange] (0,-3.8) -- (8,-3.8)
		node[midway, above, scale=0.7, fill=white, inner sep=1pt]
		{(2) BSS Configuration Response [DPP JSON Object]};
		\node[scale=0.6, gray, align=center] at (4,-4.55)
		{\textit{JSON chứa: SSID, AKM, thông số MLD, chứng chỉ mạng}};
		
		% --- Bước 3: BSS Config Result ---
		\draw[->] (8,-5.6) -- (0,-5.6)
		node[midway, above, scale=0.7, fill=white, inner sep=1pt]
		{(3) BSS Configuration Result [Success/Failure]};
		
		% --- Bước 4: Agent List ---
		\draw[->] (0,-7.0) -- (8,-7.0)
		node[midway, above, scale=0.7, fill=white, inner sep=1pt]
		{(4) Agent List Message [Cập nhật danh sách mạng]};
		\node[scale=0.6, gray, align=center] at (4,-7.75)
		{\textit{Controller cập nhật cho Agent mới biết topo toàn bộ mạng hiện tại}};
	\end{tikzpicture}
\end{center}

\textbf{Giải thích luồng bản tin:}
\begin{enumerate}
	\item Agent gửi \textit{BSS Configuration Request} liệt kê năng lực chi tiết của
	từng radio theo thứ tự HT $\to$ VHT $\to$ HE $\to$ EHT, kèm RUID để
	Controller nhận diện radio tương ứng.
	\item Controller phân tích và trả về \textit{BSS Configuration Response} chứa
	\textit{DPP Config Object} định dạng JSON --- một đối tượng cho mỗi BSS cần
	cấu hình, bao gồm tên mạng, cơ chế xác thực, tham số Wi-Fi 7 và thông tin MLD.
	\item Agent áp dụng cấu hình nhận được và gửi \textit{BSS Configuration Result}
	xác nhận thành công hoặc báo lỗi nếu có vấn đề về tương thích.
	\item Controller gửi \textit{Agent List Message} giúp Agent mới nhận biết đầy đủ
	topo mạng --- danh sách các Agent hiện có, BSSID và vai trò của từng node ---
	phục vụ tối ưu hóa Backhaul và chuyển vùng liền mạch.
\end{enumerate}


% ============================================================
\section{Đồng Bộ Chính Sách Và MLO}
% ============================================================

% ------------------------------------------------------------
\subsection{Cấu Hình Chính Sách (Policy Configuration)}
% ------------------------------------------------------------

Sau khi AP được cấu hình, Controller phân phối các \textit{chính sách vận hành} tới
từng Agent thông qua bản tin \textbf{Multi-AP Policy Config Request}~[15]. Các Agent
\textbf{phải gửi 1905 Ack trong vòng 1 giây} sau khi nhận bản tin này, trước khi bắt
đầu thực thi chính sách --- đảm bảo Controller biết chính sách đã được tiếp nhận.

Ba nhóm chính sách chính được đồng bộ:
\begin{itemize}
	\item \textbf{Steering Policy} (\textit{Chính sách chuyển vùng}):
	Định nghĩa ngưỡng cường độ tín hiệu RCPI (\textit{Received Channel Power
		Indicator}) mà dưới ngưỡng đó Agent chủ động đề xuất chuyển vùng thiết bị
	sang một điểm truy cập tốt hơn. Ngoài ra, Controller có thể chỉ định danh
	sách thiết bị \textit{không} được tự động chuyển vùng (ví dụ: thiết bị IoT
	cố định)~[15, 16].
	\item \textbf{Metric Reporting Policy} (\textit{Chính sách báo cáo số liệu}):
	Xác định chu kỳ và loại số liệu mà Agent cần định kỳ gửi lên Controller,
	bao gồm chất lượng liên kết (RSSI, throughput, PER) và tình trạng kênh
	vô tuyến~[17, 18].
	\item \textbf{Traffic Separation Policy} (\textit{Chính sách phân tách lưu lượng}):
	Ánh xạ từng SSID vào một VLAN ID cụ thể, đảm bảo lưu lượng các mạng khách
	nhau (ví dụ: mạng gia đình và mạng IoT) được cách ly hoàn toàn về mặt
	Layer 2~[19, 20].
\end{itemize}


% ------------------------------------------------------------
\subsection{Cấu Hình Wi-Fi 7 --- Đa Liên Kết (MLO)}
% ------------------------------------------------------------

Wi-Fi 7 (IEEE 802.11be) giới thiệu kiến trúc \textbf{Multi-Link Operation (MLO)}, cho
phép một thiết bị duy trì đồng thời nhiều liên kết trên các băng tần khác nhau (2.4, 5
và 6\,GHz) để tăng thông lượng và giảm độ trễ. Trong mạng Multi-AP, Controller quản lý
các thực thể MLD thông qua bản tin \textbf{AP MLD Configuration Request}~[21, 22].

Bản tin này cho phép Controller thực hiện:
\begin{itemize}
	\item \textbf{Thêm Affiliated AP:} Đưa một radio mới vào thực thể MLD hiện có,
	mở rộng số liên kết đồng thời có thể cung cấp cho thiết bị người dùng.
	\item \textbf{Xóa Affiliated AP:} Loại bỏ một radio khỏi MLD (ví dụ: khi radio
	đó cần bảo trì hoặc bị quá tải).
	\item \textbf{Cập nhật tham số MLD:} Điều chỉnh chính sách phân phối lưu lượng
	giữa các liên kết trong cùng một MLD.
\end{itemize}

Sau khi nhận bản tin, Agent áp dụng cấu hình và phản hồi trạng thái qua bản tin
\textbf{AP MLD Configuration Response}~[22], xác nhận thành công hoặc báo lỗi
nếu cấu hình không tương thích với phần cứng.


% ============================================================
\section{Báo Cáo Trạng Thái Vận Hành}
% ============================================================

Sau khi áp dụng cấu hình, Agent không chỉ thực thi mà còn phải \textbf{xác nhận
	ngược lại} với Controller về trạng thái thực tế đang vận hành. Điều này quan trọng
vì cấu hình được yêu cầu có thể không khớp hoàn toàn với phần cứng thực tế
(ví dụ: không hỗ trợ một kênh cụ thể).

Agent báo cáo thông qua \textbf{1905 Topology Response} mang \textbf{AP Operational
	BSS TLV}~[23--25], trong đó liệt kê:
\begin{itemize}
	\item BSSID thực tế đang hoạt động trên từng radio (xác định qua RUID).
	\item Kênh vô tuyến đang sử dụng và băng thông kênh thực tế.
	\item Vai trò hiện tại của từng BSS (Fronthaul / Backhaul).
\end{itemize}

Controller sử dụng thông tin này để cập nhật bản đồ topo mạng và phát hiện sớm
bất kỳ sai lệch nào giữa cấu hình được yêu cầu và trạng thái thực tế triển khai.


% ============================================================
\section{Các Điểm Kỹ thuật Then Chốt}
% ============================================================

\begin{itemize}
	\item \textbf{Định danh Radio (RUID):} Trong tất cả các bản tin đồng bộ, mỗi radio
	được xác định duy nhất bởi \textit{Radio Unique Identifier} (RUID) --- giá trị
	MAC 6 byte gán riêng cho từng radio. Controller dùng RUID để đảm bảo cấu hình
	được áp đúng radio, tránh nhầm lẫn trên các thiết bị đa băng tần.
	
	\item \textbf{Tái đồng bộ (Reconfiguration):} Khi cần thay đổi cấu hình (ví dụ: đổi
	mật khẩu Wi-Fi hoặc cập nhật chính sách), Controller buộc Agent thực hiện lại
	toàn bộ quy trình lấy cấu hình bằng cách gửi:
	\begin{itemize}
		\item \textbf{1905 AP-Autoconfiguration Renew} --- đối với luồng WSC; hoặc
		\item \textbf{Reconfiguration Trigger} --- đối với luồng DPP.
	\end{itemize}
	Agent nhận bản tin này phải ngắt cấu hình hiện tại và khởi động lại toàn bộ
	quy trình M1/M2 (WSC) hoặc BSS Configuration Request/Response (DPP).
	
	\item \textbf{Thứ tự ưu tiên và Ack:} Agent phải gửi \textit{1905 Ack} trong vòng
	\textbf{1 giây} sau khi nhận \textit{Multi-AP Policy Config Request}, trước khi
	bắt đầu thực thi chính sách mới. Cơ chế này đảm bảo Controller xác nhận rằng
	chính sách đã được Agent tiếp nhận, tránh tình trạng chính sách bị áp dụng
	không đồng nhất trên toàn mạng.
\end{itemize}